\makerubrichead{Statement of Teaching}

\begingroup
\setlength{\parindent}{0pt}
\setlength{\parskip}{6pt}
\sloppy

My teaching philosophy is centred on strong conceptual understanding, analytical reasoning, and the ability to connect mathematics with modern computational problems. I view mathematics as a foundational language for computer science, data science, artificial intelligence, optimisation, and scientific modelling; therefore, I aim to balance theoretical rigor with meaningful applications.

I have taught undergraduate and postgraduate students across BCA, MCA, B.Sc. IT, B.Sc. Data Science, M.Sc. Data Science, and M.Sc. Mathematics. Courses include \textbf{Linear Algebra, Calculus, Discrete Mathematics, Operations Research, Probability, Statistics, and Applied Mathematics}. In classroom practice I emphasise problem solving, mathematical interpretation, and links to data analysis, algorithms, optimisation, network science, and machine-learning foundations.

As \textbf{Program Head for Data Science Programs}, I also contribute to academic planning, curriculum coordination, course scheduling, faculty coordination, student mentoring, and quality processes. These responsibilities have strengthened my focus on outcome-oriented curriculum design and on maintaining a coherent mathematical foundation across computing and data-science programmes.

I encourage students to treat computational tools as extensions of mathematical thinking rather than substitutes for it. Where appropriate, I integrate Python, spreadsheets, visualisation, numerical experimentation, and reproducible workflows into teaching so that students can verify patterns, test conjectures, and connect abstract concepts with data-driven settings.

Continuous learning is an important part of my academic practice. I regularly participate in faculty development programmes, workshops, and research activities related to mathematics, graph theory, data science, artificial intelligence, and computational methods. My overall goal is to help students develop conceptual confidence, analytical maturity, and the ability to transfer mathematical reasoning into contemporary technological and research problems.

\endgroup
